\subsection{Antiferromagnetism does not remove the
obstruction}\label{proof:local_interactions:antiferromagnetism-does-not-remove-the-obstruction}

\textbf{Proposition L3 (gauge transfer).} Suppose
\(\sigma_i\in\{-1,1\}\) and \(J'_{ij}=\sigma _i\sigma _jJ_{ij}\). The
substitution \(Y_i=\sigma _iX_i\) takes the Gibbs law for J' to the
Gibbs law for J, preserves global polynomial boundedness and degree, and
preserves individual squared influences after the same substitution.

\textbf{Proof.} \(J'_{ij}X_iX_j=J_{ij}Y_iY_j\), and flipping
\(X_i\) flips \(Y_i\). \qed

On a bipartite graph, assign opposite \(\sigma\) values to its two
parts. Ferromagnetic \(J_{ij}=1/d_{0}\) becomes the antiferromagnetic
\(J'_{ij}=-1/d_{0}\) on every edge. Thus the same bounded-degree example
fails AA for the fixed antiferromagnetic term as well, witnessed by

\[
M_{\rm stag}(x)=n^{-1}\sum_i\sigma_ix_i.
\]

In particular the sign of the interaction by itself is not a universal
protection against AA failure. These bipartite antiferromagnets are
unfrustrated, so frustration is not necessary either.
